\subsection{From time-frequency analysis to free boundary problems}

Having translated, in the previous subsection, the time--frequency
localization problem into a question about the vanishing of a Fourier
transform along the two complex lines $\eta=\pm i\xi$, we now bridge the gap
between this Fourier-analytic statement and the free boundary formulation
that will dominate the remainder of the paper. The bridge is provided by a
divisibility lemma of Ehrenpreis--Malgrange flavour: any compactly supported
distribution whose Fourier transform vanishes on the characteristic variety
$\{\eta=\pm i\xi\}$ of the Laplacian must be the Laplacian of a compactly
supported $L^2$ function which, in our setting, is moreover continuous off a
finite set of points.

In order to state the lemma cleanly, we make precise the class of
distributions with which we work. We say that a compactly supported
distribution $\omega$ is a \emph{singular-continuous free measure} if it
defines a (finite, signed) Radon measure on $\R^2$ whose Lebesgue
decomposition with respect to planar Lebesgue measure $\mathrm{d}x\,
\mathrm{d}y$ takes the form
\[
\omega = \omega_{\mathrm{ac}} + \omega_{\mathrm{pp}}
       = f\,\mathrm{d}x\,\mathrm{d}y + \sum_{j=1}^N c_j\,\delta_{p_j},
\]
with \(f\in L^\infty(\R^2)\) compactly supported, \(c_j\in \R\setminus\{0\}\), and
$p_1,\ldots,p_N\in \R^2$ pairwise distinct. In other words, the continuous
singular part $\omega_{\mathrm{cs}}$ in the Lebesgue--Radon--Nikodym
decomposition $\omega=\omega_{\mathrm{ac}}+\omega_{\mathrm{cs}}+
\omega_{\mathrm{pp}}$ is required to vanish identically, and the pure
point part is finite. The terminology \emph{free measure} is suggested by
the role of $\omega$ as the right-hand side of a free boundary problem
arising in the next subsection.

\begin{lemma}\label{lemma:fourier-to-fboundary}
Let $\omega\in \mathcal{S}'(\R^2)$ be a compactly supported singular-continuous
free distribution such that
\[
\widehat{\omega}(w,\pm iw)=0,\qquad \forall w\in\C.
\]
Then there exist a finite set $P\subset \R^2$ and a compactly supported
function $u\in L^2(\R^2)$ such that
\[
\Delta u=\omega
\]
in the sense of distributions, and
\[
u\in C(\R^2\setminus P).
\]
\end{lemma}

\begin{proof}
By the Paley--Wiener--Schwartz theorem (see, e.g.,
\cite{Hormander,Folland}), the Fourier
% TODO(bib): the keys \texttt{Hormander} (Ark. Mat. 1991) and
% \texttt{Folland} (Harmonic Analysis in Phase Space) do not contain
% Paley--Wiener--Schwartz at the cited locations; a textbook entry
% (Hormander Vol.~I or Rudin's Functional Analysis) should be added.
transform $\widehat\omega(\xi,\eta)$ extends to an entire function on $\C^2$
of exponential type, satisfying the bound
\[
|\widehat\omega(\xi,\eta)|\le C(1+|\xi|+|\eta|)^{N}\,e^{R(|\Im\xi|+|\Im\eta|)}
\]
for some $C,R>0$ and some integer $N\ge 0$ depending on the order of
$\omega$, where $R$ is the radius of any disk containing
$\operatorname{supp}\omega$. In particular, $\widehat\omega$ admits an
everywhere convergent power series expansion
\[
\widehat{\omega}(\xi,\eta)=\sum_{k,l\ge 0} a_{k,l}\,\xi^k\eta^l,
\qquad (\xi,\eta)\in\C^2.
\]
Using the vanishing hypothesis on the complex lines $\eta=\pm i\xi$, we may
group terms by total homogeneity: setting $\eta=\pm iw$, $\xi=w$, and
collecting powers of $w$,
\[
0=\widehat{\omega}(w,\pm iw)
   =\sum_{n\ge 0}\Bigl(\sum_{k=0}^n a_{k,n-k}(\pm i)^{\,n-k}\Bigr)w^n,
\qquad \forall w\in\C.
\]
Since this holds for all $w\in\C$, comparing coefficients yields
\[
\sum_{k=0}^n a_{k,n-k}(\pm i)^{\,n-k}=0,\qquad \forall n\ge 0.
\]
For each $n\ge 0$, define the auxiliary one-variable polynomial
\[
p_n(t):=\sum_{k=0}^n a_{k,n-k}\, t^{\,n-k}.
\]
The two scalar equations above precisely say that $p_n(\pm i)=0$, so
$(t-i)(t+i)=t^2+1$ divides $p_n(t)$. Hence we may write
\[
p_n(t)=(t^2+1)\,q_n(t),
\]
with $q_n$ a polynomial of degree at most $n-2$ (and $q_n\equiv 0$ when
$n<2$).

We now reassemble the homogeneous pieces. Let
\[
H_n(\xi,\eta):=\sum_{k=0}^n a_{k,n-k}\,\xi^k\eta^{\,n-k},
\qquad n\ge 0,
\]
so that $\widehat\omega=\sum_{n\ge 0}H_n$ is the decomposition of
$\widehat\omega$ into homogeneous components (of degree $n$). The argument
above shows that each $H_n$ vanishes on the two complex lines
$\eta=\pm i\xi$, hence is divisible, as a homogeneous polynomial in two
variables, by $\xi^2+\eta^2$. Concretely,
\[
H_n(\xi,\eta)=(\xi^2+\eta^2)\,G_{n-2}(\xi,\eta)
\]
for a (uniquely determined) homogeneous polynomial $G_{n-2}$ of degree
$n-2$, with $G_{n-2}\equiv 0$ for $n<2$. Summing the homogeneous pieces, the
formal series
\[
F(\xi,\eta):=\sum_{n\ge 2} G_{n-2}(\xi,\eta)
\]
converges absolutely on every polydisc on which $\widehat\omega$ does, and
defines an entire function on $\C^2$ satisfying
\[
\widehat{\omega}(\xi,\eta)=(\xi^2+\eta^2)\,F(\xi,\eta).
\]
Since $\widehat\omega$ has a zero of order at least two at the origin
(coming from the vanishing of $H_0$ and $H_1$, which is forced by
$p_0(\pm i)=p_1(\pm i)=0$), this factorization is consistent at $0$ and $F$
is entire there as well.

We next quantify the size of $F$ on $\R^2$. Since $\omega$ is a finite
(signed) Radon measure with compact support, its Fourier transform is in
fact bounded on $\R^2$:
\[
|\widehat{\omega}(\xi,\eta)|\le \|\omega\|_{\mathrm{TV}},
\qquad (\xi,\eta)\in\R^2,
\]
where $\|\omega\|_{\mathrm{TV}}$ is the total variation; equivalently,
the Paley--Wiener bound holds with $N=0$. On $\R^2$, $\xi^2+\eta^2=
|\xi|^2+|\eta|^2$, so for $|\xi|^2+|\eta|^2\ge 1$ we have $\xi^2+\eta^2
\ge \tfrac{1}{2}(1+|\xi|^2+|\eta|^2)$, whence
\[
|F(\xi,\eta)|=\frac{|\widehat\omega(\xi,\eta)|}{\xi^2+\eta^2}
\le \frac{2\|\omega\|_{\mathrm{TV}}}{1+|\xi|^2+|\eta|^2},
\qquad |\xi|^2+|\eta|^2\ge 1.
\]
Near the origin, $F$ is entire (as established above) and hence bounded;
absorbing this region into the constants we obtain
\[
|F(\xi,\eta)|\le \frac{C''}{1+|\xi|^2+|\eta|^2},
\qquad (\xi,\eta)\in\R^2,
\]
for some $C''>0$. Therefore $F|_{\R^2}\in L^2(\R^2)\cap L^1(\R^2)$. On
$\C^2$, $F$ is entire (by the formal-series construction above) and
satisfies an exponential-type bound of the same rate $R$ as
$\widehat\omega$. Indeed, on the open set
$\Omega:=\{(\xi,\eta)\in\C^2:|\xi^2+\eta^2|\ge 1\}$ one has the pointwise
estimate
\[
|F(\xi,\eta)|=\frac{|\widehat\omega(\xi,\eta)|}{|\xi^2+\eta^2|}
\le C(1+|\xi|+|\eta|)^{N}\,e^{R(|\Im\xi|+|\Im\eta|)},
\qquad (\xi,\eta)\in\Omega,
\]
by the Paley--Wiener bound on $\widehat\omega$. On the complement
$\Omega^c=\{|\xi^2+\eta^2|<1\}$, the function $F$ is entire and the
maximum modulus principle, applied on each complex line through a given
point in $\Omega^c$ in a direction transverse to the variety
$\{\xi^2+\eta^2=0\}$, bounds $|F|$ at points of $\Omega^c$ by a uniform
constant times the supremum of $|F|$ on the boundary of a fixed thin
tubular neighbourhood of that variety, where the previous estimate
applies. Combining these two regimes and absorbing the resulting
bookkeeping constant yields, after enlarging $C'$ if necessary,
\[
|F(\xi,\eta)|\le C'\,(1+|\xi|+|\eta|)^{N}\,
e^{R(|\Im \xi|+|\Im \eta|)},
\qquad (\xi,\eta)\in\C^2,
\]
which is the Paley--Wiener--Schwartz bound for a compactly supported
distribution of order $\le N$ in $\overline{B(0,R)}$.

By the Paley--Wiener--Schwartz theorem applied in the reverse direction
(see again \cite{Hormander}), the entire function $F$
satisfies an exponential-type bound on $\C^2$ with the same exponential
rate $R$ as $\widehat\omega$, and integrable decay along $\R^2$, so it is
the Fourier--Laplace transform of a distribution of order $\le N$
compactly supported in the closed disk $\overline{B(0,R)}$. The $L^2$
bound on the real locus shows that this distribution is, in fact, an
$L^2(\R^2)$ function. Let $u\in L^2(\R^2)$ denote this compactly supported
representative, characterised by $\widehat u = F$ on $\R^2$ and
$\operatorname{supp} u\subset \overline{B(0,R)}$.

Computing distributional Laplacians on the Fourier side, with the
convention $\widehat g(\xi,\eta)=\int g(x,y)e^{-2\pi i(x\xi+y\eta)}\,
\mathrm{d}x\,\mathrm{d}y$, we have
\[
\widehat{\Delta u}(\xi,\eta) = -4\pi^2(\xi^2+\eta^2)\,\widehat{u}(\xi,\eta)
= -4\pi^2(\xi^2+\eta^2)\,F(\xi,\eta) = -4\pi^2\,\widehat\omega(\xi,\eta);
\]
with any other standard normalisation, the same identity holds up to a
nonzero real multiplicative factor. In all cases, $\Delta u$ is a constant
nonzero real multiple of $\omega$. Replacing $u$ by $\lambda u$ for the
appropriate (nonzero, real) $\lambda$ -- this is the harmless rescaling
referred to above -- we may, and do, assume that
\[
\Delta u = \omega
\]
in the sense of distributions. The rescaling preserves all the previously
established properties of $u$: compact support, membership in $L^2(\R^2)$,
and the Fourier identity $\widehat u = \lambda F$.

To pin down $u$ pointwise, observe that since $\omega$ is a compactly
supported Radon measure on $\R^2$, the convolution
\[
v(x,y) := \frac{1}{2\pi}(\log|\cdot|*\omega)(x,y)
\]
is a well-defined locally integrable function on $\R^2$ which solves
$\Delta v = \omega$ in the sense of distributions; here $G(z)=\log|z|$ is
(up to the standard normalisation $\tfrac{1}{2\pi}$) the fundamental solution
of the Laplacian in $\R^2$, in the sense that
$\Delta\bigl(\tfrac{1}{2\pi}\log|\cdot|\bigr)=\delta_0$, see, e.g.,
\cite{Hormander} for the underlying distribution theory and
\cite{Aharonov-Schiffer-Zalcman} for related representations in the planar
quadrature-domain context.
% TODO(bib): \texttt{Garnett} (Bounded Analytic Functions, GTM 236) does
% not contain the previously-cited Ch.~II Thm~1.3 on continuity of
% logarithmic potentials of bounded densities; consider adding Ransford
% (Potential Theory in the Complex Plane) or Landkof (Foundations of
% Modern Potential Theory) to the bibliography.
The difference $u-v$ is then a tempered
distribution which is harmonic on $\R^2$ (since
$\Delta(u-v)=\omega-\omega=0$), and hence, being harmonic, real-analytic.
Outside the union of the supports of $u$ and $\omega$, the function $v$ is
in fact real-analytic and grows at most logarithmically at infinity, with
$v(x)=\frac{M}{2\pi}\log|x|+O(|x|^{-1})$ as $|x|\to\infty$ where
$M=\omega(\R^2)$, while $u$ vanishes there because of its compact support.
Hence $u-v$ is harmonic on $\R^2$ and grows at most logarithmically at
infinity; a standard refinement of the harmonic Liouville theorem (any
harmonic function on $\R^2$ growing slower than a polynomial of degree $1$
is constant) then forces $u-v$ to be a constant, and this constant must
vanish to be consistent with $u\in L^2(\R^2)$ and the local integrability
of $v$. Equivalently, and more transparently: by the compact support of
$\omega$ together with $\widehat u = F$ on $\R^2$, the function $u$ is
uniquely determined among compactly supported $L^2$ solutions of
$\Delta u = \omega$, since any two such would differ by an entire $L^2$
function which is harmonic, hence vanishes. Thus
\[
u(x,y)=\frac{1}{2\pi}\int_{\R^2}\log|(x,y)-(x',y')|\,\mathrm{d}\omega(x',y'),
\]
the logarithmic potential of $\omega$.

We finally analyse the regularity of $u$. Decomposing
$\omega = f\,\mathrm{d}x\,\mathrm{d}y+\sum_{j=1}^N c_j\,\delta_{p_j}$ as in
the definition of singular-continuous free measure, we obtain
correspondingly
\[
u(x,y)=\frac{1}{2\pi}\int_{\R^2}\log|(x,y)-(x',y')|\,f(x',y')\,
\mathrm{d}x'\mathrm{d}y'
+\frac{1}{2\pi}\sum_{j=1}^N c_j\,\log|(x,y)-p_j|.
\]
The first summand is the logarithmic potential of the bounded, compactly
supported density \(f\in L^\infty(\R^2)\) and is therefore not merely
continuous but in fact $C^{1,\alpha}(\R^2)$ for every $\alpha\in(0,1)$
(by the standard Calder\'on--Zygmund / Schauder estimates applied to
$\Delta=\partial_x^2+\partial_y^2$ with right-hand side in $L^\infty$
of compact support, see, e.g., the Newtonian-potential discussion in
\cite{Hormander}).
The continuity used here is in fact elementary: by dominated convergence,
$z\mapsto \int \log|z-w|\,f(w)\,\mathrm{d}A(w)$ is continuous since
$\log|\cdot-w|$ is locally integrable in $w$ uniformly on compacta in $z$,
and $f$ is bounded with compact support. It is here that the absence of a continuous
singular part in $\omega$ is essential: a generic continuous singular
measure (e.g.\ the Hausdorff measure on a Cantor-type set of positive
$1$-capacity) gives rise to a logarithmic potential which need not be
continuous, only quasi-continuous, and the conclusion of the lemma would
fail in that generality. Each term $\log|(x,y)-p_j|$ is continuous on
$\R^2\setminus\{p_j\}$, with a logarithmic singularity at $p_j$. Setting
$P:=\{p_1,\ldots,p_N\}$, we conclude that $u\in C(\R^2\setminus P)$, as
claimed. The fact that $u$ is compactly supported has already been
established in the construction above and is consistent with the
representation as a logarithmic potential, which differs from a compactly
supported function only by a globally harmonic correction that, as
explained, must vanish identically.
\end{proof}

We postpone the proof of Theorem~\ref{thm:eig-loc-U} until after the
perturbative inverse theorem. The same inverse/free-boundary mechanism used
to construct domains for prescribed eigenfunctions also recovers the
classical Abreu--D\"orfler disk theorem in the simply connected setting; see
Section~\ref{sec:classical-rigidity} below.


% ==============================================================
% Section 3
% ==============================================================
